\documentclass[11pt, a4paper]{article}
\usepackage[utf8]{inputenc}
\usepackage[T1]{fontenc}
\usepackage[brazil]{babel}
\usepackage{geometry}
\geometry{a4paper, margin=2cm}
\usepackage{graphicx}
\usepackage{tabularx}
\usepackage{colortbl}
\usepackage{xcolor}
\usepackage{booktabs}
\usepackage{enumitem}
\usepackage{hyperref}
\usepackage{helvet}
\renewcommand{\familydefault}{\sfdefault} % Arial-like font

\definecolor{headerblue}{RGB}{20, 50, 100} % Dark blue for headers
\definecolor{rowgray}{RGB}{240, 240, 240} % Light gray for table rows

\begin{document}

% Cabeçalho
\begin{center}
\renewcommand{\arraystretch}{1.5}
\begin{tabularx}{\textwidth}{|X|X|l|l|}
\hline
\rowcolor{headerblue}
\multicolumn{4}{|c|}{\textcolor{white}{\textbf{\Large DECLARAÇÃO DE ESCOPO}}} \\
\hline
Nome do Projeto: & \multicolumn{3}{l|}{Projeto Conecta - Plataforma de Conexão Local (Android \& Web)} \\
\hline
Gerente do Projeto: & \multicolumn{3}{l|}{João Bruno Faria Rezende} \\
\hline
Data: & 28/08/2026 & Nº Versão do Documento: & 2.0 \\
\hline
\end{tabularx}
\end{center}

\vspace{0.5cm}

\section*{\textcolor{headerblue}{1. Título e Descrição do projeto}}
\textbf{Projeto Conecta}\\
O projeto Conecta consiste no desenvolvimento de uma plataforma digital que integra usuários comuns, lojistas e prestadores de serviços em um único aplicativo móvel e web, facilitando a busca, contratação, avaliação e gerenciamento de produtos e serviços locais e regionais. Desenvolvido no contexto da F.A.S.T (Fábrica de Soluções Tecnológicas) em parceria com o IFMG – Campus Bambuí.

O projeto implantará uma plataforma desenvolvida em Flutter com backend em Firebase. A solução abrangerá controle de acesso avançado (Lazy Authentication para visitantes, Double Opt-in para cadastros), perfis distintos (Comum, Lojista, Prestador e Administrador), vitrine de produtos e serviços, agendamento de serviços, carrinho de compras, e painel administrativo para validação de documentos de lojas e prestadores. O lançamento inicial atenderá usuários a nível regional, sendo exclusivo para dispositivos Android (versão 8.0+) e navegadores Web.

O escopo do projeto compreende o levantamento (apoiado por orientadores do IFMG), design de interfaces, desenvolvimento, testes, publicação na Google Play Store e versão Web, suporte inicial e encerramento.

\section*{\textcolor{headerblue}{2. Patrocinador (sponsor)}}
F.A.S.T (Fábrica de Soluções Tecnológicas) e Instituto Federal de Minas Gerais (IFMG) – Campus Bambuí. Responsabilidades: prover orientação acadêmica (Me. Cláudio Ribeiro de Sousa e Dr. Eduardo Cardoso Melo), aprovar a linha de base de escopo e formalizar o aceite das entregas.

\section*{\textcolor{headerblue}{3. Gerente do Projeto e Autoridade}}
João Bruno Faria Rezende, Gerente de Projetos. Possui autoridade para planejar e coordenar a equipe de desenvolvimento, priorizar requisitos, convocar validações e assegurar o cumprimento dos objetivos técnicos.

\section*{\textcolor{headerblue}{4. Usuários}}
\begin{enumerate}[label=\arabic*.]
    \item \textbf{Usuário Comum (Cliente):} pesquisa produtos e serviços, adiciona ao carrinho, favorita, agenda serviços, acompanha pedidos e avalia após o consumo real.
    \item \textbf{Lojista:} gerencia catálogo de produtos, recebe pedidos, confirma estoque manualmente e atende clientes. Pode ser empresa formal (com CNPJ) ou autônomo (sem CNPJ).
    \item \textbf{Prestador de serviço:} disponibiliza portfólio, define áreas de atuação, gerencia agenda de serviços e status de disponibilidade (Online/Offline).
    \item \textbf{Administrador:} acessa painel exclusivo para aprovar ou rejeitar (com justificativa) cadastros e documentos de Lojistas e Prestadores, além de moderar avaliações.
    \item \textbf{Visitante:} navega pelo catálogo e visualiza itens sem autenticação (Lazy Authentication). O login é exigido apenas ao tentar comprar, favoritar ou interagir.
\end{enumerate}

\section*{\textcolor{headerblue}{5. Equipe}}
\begin{center}
\renewcommand{\arraystretch}{1.3}
\begin{tabularx}{\textwidth}{|l|l|X|}
\hline
\rowcolor{headerblue}
\textcolor{white}{\textbf{Papel}} & \textcolor{white}{\textbf{Nome}} & \textcolor{white}{\textbf{Responsabilidade principal}} \\
\hline
Gerente de Projetos & João Bruno Faria Rezende & Governança, prazos, documentação, gestão do backlog e coordenação \\
\hline
Desenvolvedor & Gabriel Henrique Silva Duque & Implementação frontend/backend (Flutter/Firebase), integrações e UI \\
\hline
Desenvolvedor & Rafaek Oliveira Gonçalves & Implementação frontend/backend (Flutter/Firebase), testes e regras de negócio \\
\hline
\end{tabularx}
\end{center}

\section*{\textcolor{headerblue}{6. Objetivos}}
\begin{enumerate}[label=\arabic*.]
    \item Lançar a plataforma (Android e Web) até o final do cronograma acadêmico.
    \item Fornecer uma interface unificada que reduza em 80\% o tempo gasto por consumidores na busca de fornecedores locais.
    \item Garantir um tempo de resposta inferior a 3 segundos para interações na interface, com leitura de banco de dados abaixo de 1 segundo.
    \item Implementar um processo de moderação e validação de 100\% dos lojistas e prestadores de serviços via painel administrativo.
\end{enumerate}

\section*{\textcolor{headerblue}{7. Justificativa}}
O comércio local pulverizado carece de ferramentas digitais robustas, seguras e baratas. O Conecta visa promover a transformação digital e o fortalecimento do comércio de Bambuí-MG e região. A aplicação de Lazy Authentication garante baixa fricção de entrada, enquanto a validação de documentos de lojistas proporciona segurança aos consumidores.

\section*{\textcolor{headerblue}{8. Fatores de Sucesso}}
\begin{enumerate}[label=\arabic*.]
    \item Aceitação positiva na validação do protótipo e navegação intuitiva (máximo de 3 cliques para ações principais).
    \item Estabilidade no banco de dados, suportando 500 usuários simultâneos sem perda perceptível de performance.
    \item Segurança contra fraudes através da aprovação estrita de documentos pelo perfil Administrador.
    \item Avaliações reais e orgânicas, restritas a usuários com histórico de consumo confirmado na plataforma.
\end{enumerate}

\section*{\textcolor{headerblue}{9. Restrições}}
\begin{enumerate}[label=\arabic*.]
    \item \textbf{Plataforma:} O aplicativo não será desenvolvido ou publicado para o ecossistema iOS (Apple), limitando-se ao Android e Web.
    \item \textbf{Pagamentos:} A plataforma \textbf{não realiza transações financeiras nem processa pagamentos}. Acordos de pagamento e entrega são feitos diretamente entre vendedor e cliente (ex: pagamento na entrega).
    \item \textbf{Equipe:} O desenvolvimento será realizado estritamente pelos três membros definidos, limitando a capacidade produtiva diária.
\end{enumerate}

\section*{\textcolor{headerblue}{10. Premissas}}
\begin{enumerate}[label=\arabic*.]
    \item A F.A.S.T/IFMG disponibilizará os recursos básicos de orientação acadêmica.
    \item O Firebase Spark (gratuito) será suficiente para a fase inicial de operação e testes.
    \item Lojistas e Prestadores serão responsáveis por alimentar seus catálogos e manter o status de disponibilidade (Online/Offline) atualizado.
\end{enumerate}

\section*{\textcolor{headerblue}{11. Requisitos funcionais}}
\begin{enumerate}[label=RF-\arabic*.]
    \item \textbf{Acesso Visitante:} O sistema deve usar \textit{Lazy Authentication}, permitindo visualização de lojas sem login, armazenando a intenção de ação para execução após a autenticação.
    \item \textbf{Cadastros Específicos:} O sistema deve coletar CPF/CNPJ, CNAE, área de atuação e documentos de certificação, com validação de formato e máscara em tempo real.
    \item \textbf{Painel Admin:} Administradores devem ter uma fila de aprovação de cadastros, podendo Aprovar ou Rejeitar (com envio de feedback do motivo) os lojistas/prestadores.
    \item \textbf{Vitrine e Agendamento:} Clientes podem adicionar produtos de lojistas à sacola ou selecionar horários na agenda de prestadores de serviço.
    \item \textbf{Confirmação de Estoque:} O lojista deve aprovar o pedido manualmente, validando o estoque real antes da continuidade.
    \item \textbf{Avaliações Concretas:} Avaliações em estrelas (1 a 5) só podem ser submetidas após a confirmação de conclusão do pedido/serviço.
    \item \textbf{Notificações:} Enviar alertas com sons distintos baseados no perfil do usuário (ex: toque de nova venda para lojistas).
\end{enumerate}

\section*{\textcolor{headerblue}{12. Requisitos não funcionais}}
\begin{enumerate}[label=RNF-\arabic*.]
    \item \textbf{Desempenho:} Resposta à interação do usuário em até 3 segundos e leitura do banco de dados em até 1 segundo. Suporte inicial a 500 usuários simultâneos.
    \item \textbf{Usabilidade:} Navegação com no máximo 3 cliques para ações principais. Correções de arquitetura de rotas para não permitir empilhamento infinito de páginas (Home Bars).
    \item \textbf{Segurança:} Dados sensíveis (senhas) criptografados. Uso de HTTPS e proteção contra injeções.
    \item \textbf{Legalidade:} Cumprimento da LGPD, requerendo aceite explícito dos Termos de Uso e clareza sobre a isenção do aplicativo em transações financeiras.
    \item \textbf{Compatibilidade:} Funcionamento em Android 8.0+ e navegadores web modernos.
\end{enumerate}

\section*{\textcolor{headerblue}{13. Exclusões Específicas}}
\begin{enumerate}[label=\arabic*.]
    \item \textbf{Ecossistema Apple:} Compilação, testes ou publicação na Apple App Store (iOS).
    \item \textbf{Meios de Pagamento:} Gateway de cartão de crédito ou processamento nativo de transações financeiras de qualquer natureza.
    \item \textbf{Logística Própria:} O sistema não fornece frota nem cálculo avançado de correios, sendo a entrega combinada entre as partes.
\end{enumerate}

\section*{\textcolor{headerblue}{14. Entregas Principais}}
\begin{center}
\renewcommand{\arraystretch}{1.3}
\begin{tabularx}{\textwidth}{|l|l|X|}
\hline
\rowcolor{headerblue}
\textcolor{white}{\textbf{Código}} & \textcolor{white}{\textbf{Entrega}} & \textcolor{white}{\textbf{Conteúdo verificável}} \\
\hline
1.1 & Levantamento e UX & Protótipo do Figma validado com orientadores e stakeholders \\
\hline
1.2 & Módulo de Autenticação & Double Opt-in, Lazy Authentication, máscaras de CPF/CNPJ e validações \\
\hline
1.3 & Módulo Administrativo & Painel de aprovação/rejeição de documentos e cadastros pendentes \\
\hline
1.4 & Visão do Consumidor & Vitrines, carrinho de compras, agendamento de serviços e avaliações \\
\hline
1.5 & Painéis Comerciais & Dashboards de lojistas e prestadores para aceitar pedidos e mudar status \\
\hline
1.6 & Publicação & Versão Web online e APK gerado para Android (Play Store) \\
\hline
\end{tabularx}
\end{center}

\section*{\textcolor{headerblue}{15. Marcos Principais}}
\begin{center}
\renewcommand{\arraystretch}{1.3}
\begin{tabularx}{\textwidth}{|l|l|X|}
\hline
\rowcolor{headerblue}
\textcolor{white}{\textbf{Data}} & \textcolor{white}{\textbf{Marco}} & \textcolor{white}{\textbf{Condição de conclusão}} \\
\hline
Etapa 1 & Prototipação e Requisitos & Documento de Requisitos e Figma concluídos \\
\hline
Etapa 2 & Estruturação Backend & Firebase modelado (Firestore, Auth, Storage) \\
\hline
Etapa 3 & Frontend Android/Web & Telas responsivas conectadas às regras de negócio \\
\hline
Etapa 4 & Painel Administrativo & Fluxo de validação de documentos concluído \\
\hline
Etapa 5 & Testes e Homologação & Tratamento de falhas e performance (ex: correção do Menu Lateral) \\
\hline
Etapa 6 & Lançamento Final & Defesa do projeto e publicação do APK/Web \\
\hline
\end{tabularx}
\end{center}

\section*{\textcolor{headerblue}{16. Critérios de Aceitação do Projeto}}
\begin{enumerate}[label=\arabic*.]
    \item Plataforma instalável em dispositivos Android e acessível via navegador Web.
    \item Painel administrativo plenamente capaz de bloquear acesso comercial de lojistas/prestadores até a validação de seus documentos.
    \item App isento de transações financeiras, respeitando estritamente os Termos de Uso.
    \item Sistema tolerante a 500 usuários simultâneos mantendo tempo de resposta ágil.
    \item Código-fonte estruturado no GitHub, com relatórios técnicos submetidos aos orientadores do IFMG.
\end{enumerate}

\vspace{1cm}
\noindent
\textbf{Aprovado por:} Orientadores do Projeto Conecta (F.A.S.T/IFMG) \\
\textbf{Assinatura:} \hrulefill \quad \textbf{Data:} \_\_\_/\_\_\_/\_\_\_\_\_

\newpage

\section*{\textcolor{headerblue}{ANEXO A - EAP COMPLETA}}
\textbf{Projeto Conecta - Plataforma de Conexão Local}

\begin{center}
\renewcommand{\arraystretch}{1.3}
\begin{tabularx}{\textwidth}{|c|X|X|c|}
\hline
\rowcolor{headerblue}
\textcolor{white}{\textbf{Código}} & \textcolor{white}{\textbf{Elemento da EAP}} & \textcolor{white}{\textbf{Entrega / resultado verificável}} & \textcolor{white}{\textbf{Responsável}} \\
\hline
\textbf{1.0} & \textbf{Projeto Conecta} & \textbf{Aplicativo Android e Web para comércio local.} & \textbf{Equipe} \\
\hline
\rowcolor{rowgray}
1.1 & Requisitos e Design & Entregas base focadas no usuário. & João Bruno \\
\hline
1.1.1 & Documentação & Levantamento funcional e não funcional aprovado. & João Bruno \\
\hline
1.1.2 & Prototipação & Fluxos no Figma validados pelos orientadores. & João Bruno \\
\hline
\rowcolor{rowgray}
1.2 & Autenticação e Perfis & Segurança e segmentação de usuários. & Gabriel / Rafaek \\
\hline
1.2.1 & Lazy Authentication & Acesso visitante com persistência de ação após login. & Rafaek \\
\hline
1.2.2 & Cadastro Lojista/Prestador & Coleta de documentos, CNPJ/CPF com validação nativa. & Gabriel \\
\hline
\rowcolor{rowgray}
1.3 & Módulo Administrativo & Controle interno de segurança da plataforma. & Gabriel / Rafaek \\
\hline
1.3.1 & Validação de Contas & Fila de pendentes, download de documentos e aceite/recusa. & Gabriel \\
\hline
\rowcolor{rowgray}
1.4 & Visão do Consumidor & Funcionalidades de compra e agendamento. & Gabriel / Rafaek \\
\hline
1.4.1 & Catálogo e Carrinho & Seleção de produtos, variações e checkout sem pagamento. & Rafaek \\
\hline
1.4.2 & Agendamento & Seleção de dias/horários na agenda do prestador. & Gabriel \\
\hline
1.4.3 & Avaliações & Sistema de reviews restrito a histórico de compras reais. & Rafaek \\
\hline
\rowcolor{rowgray}
1.5 & Painéis Comerciais & Telas para lojistas e prestadores gerirem o negócio. & Gabriel / Rafaek \\
\hline
1.5.1 & Gestão de Estoque & Lojista recebe pedido e valida manualmente disponibilidade. & Rafaek \\
\hline
1.5.2 & Status de Serviço & Prestador define Online/Offline e responde solicitações. & Gabriel \\
\hline
\rowcolor{rowgray}
1.6 & Testes e Lançamento & Qualidade e deploy final. & Equipe \\
\hline
1.6.1 & Correção de Bugs & Estabilização de rotas (correção de vazamento de tela). & Gabriel / Rafaek \\
\hline
1.6.2 & Deploy & Build para Web e pacote APK para Android. & Equipe \\
\hline
\end{tabularx}
\end{center}

\end{document}
